\BourbakiBackSection{HISTORICAL NOTE}

(Numbers in brackets refer to the bibliography at the end of this note.)

The principal concepts and results on uniform spaces have emerged gradually from the theory of real variables, and it is only in recent years that they have been systematically studied. Cauchy, in his attempt to put the theory of series on a rigorous basis (cf.~the Historical Notes to Chapters I and IV) took as starting point a principle which he seems to have thought self-evident, namely that a necessary and sufficient condition for the convergence of a sequence \((a_{n})\) is that \(|a_{n+p} - a_{n}|\) is as small as we please whenever n is sufficiently large (see for example {[}2{]}). Together with Bolzano {[}1{]} he was undoubtedly one of the first to state this principle explicitly and to recognize its importance; hence the name of ``Cauchy sequences'' given to sequences of real numbers which satisfy the condition in question, and by extension to sequences \((x_{n})\) of points in a metric space (Chapter IX) such that the distance from \(x_{n+p}\) to \(x_{n}\) is as small as we please whenever n is large enough; and hence finally the name of ``Cauchy filter'' given to the generalization of Cauchy sequences which has been studied in this chapter.

When later the intuitive notion of real number ceased to be regarded as adequate, and attempts were made to define real numbers in terms of rational numbers in order to provide a solid basis for Analysis, it was Cauchy's principle which provided the most fertile of the definitions put forward in the second half of the nineteenth century. This definition, due to Cantor {[}3{]} (and developed along the lines of Cantor's ideas by Heine {[}5{]} and independently by Méray, amongst others) consists in making a real number correspond to every Cauchy sequence of rational numbers (``fundamental sequence'' in Cantor's terminology); the same real number corresponds to two Cauchy sequences \((a_{n})\) and \((b_{n})\) of rational numbers if and only if \(|a_{n} - b_{n}|\) tends to zero. The essential idea here is that, from a certain point of view (in fact, the point of view of the ``uniform structure'' defined in this chapter, § 1, no. 1, Example 1), the set Q of rational numbers is ``incomplete'', and that the set of real numbers is the ``complete'' set which one gets from Q by ``completion''.

On the other hand, Heine was the first to define uniform continuity for real-valued functions of one or more real variables, in work largely inspired by the ideas of Weierstrass and Cantor {[}4{]}, and he proved that every real-valued function which is continuous on a bounded closed interval of R is uniformly continuous on this interval: this is ``Heine's theorem''. By Theorem 2 of § 4 this result is related to the compactness of a bounded closed interval of R (``Borel-Lebesgue theorem'', Chapter IV, § 2, Theorem 2; cf.~the Historical Notes to Chapters I and IV), and Heine's proof of his theorem can also serve, with some modifications, to prove the Borel-Lebesgue theorem (and this has appeared as sufficient reason to some authors for calling the latter theorem the ``Heine-Borel theorem'').

These ideas became extended to more general spaces when metric spaces came to be studied, first in particular cases and then in general (cf.~Chapter IX); in a metric space a distance is given (i.e.~a real-valued function of pairs of points, satisfying certain axioms) which defines both a topology and a uniformity. Fréchet, who was the first to give a general definition of these spaces, recognized the importance of Cauchy's principle {[}6{]} and also proved for metric spaces a theorem equivalent to Theorem 3 of § 4 ({[}6{]} and {[}7{]}). Hausdorff developed considerably the theory of metric spaces in his ``Mengenlehre'' ({[}8{]}; cf.~also {[}8 a{]}) and in particular realized that one can apply Cantor's construction, described earlier, to these spaces and hence obtain a ``complete'' metric space from a ``non-complete'' metric space (i.e.~one in which Cauchy's principle is not valid).

Metric spaces are a particular type of ``uniform spaces''; the latter were defined in full generality by A. Weil at a fairly recent date {[}9{]}. Previous to this, the notions and results on ``uniform structure'' could be used only in connection with metric spaces, and this fact explains the important part played by metric spaces and metrizable spaces (and in particular by compact metrizable spaces) in many modern works on topology, in questions where distance has no real usefulness. Once one has the definition of a uniform space, there is no difficulty (especially if one also has the notion of a filter at one's disposal) in extending to these spaces almost the whole of the theory of metric spaces as given e.g.~by Hausdorff (and similarly in extending, for example, to arbitrary compact spaces the results on compact metric spaces given in Alexandroff-Hopf's Topologie {[}10{]}). This is what we have done in this chapter; in particular, the theorem on completion of uniform spaces (§ 3, Theorem 3) is no more than a transposition, without any essential modifications, of Cantor's construction of the real numbers.

